%% notes-tex/019-simb-multimodal/sections/10-launch.tex
%% SOURCE: experiments/019-simb-multimodal/results/launch_plan_evidence.json, generated by
%% experiments/019-simb-multimodal/scripts/analyze_launch_plan_evidence.py from
%% round_leaderboards.csv, the packing benchmark logs under results/packing_benchmark/
%% (raw output of scripts/gh_packing_benchmark.sh), and the 1445-last.ckpt state dict.
%% The allocation is a proposal; every number feeding it is measured and sourced.

\section{The launch, 2026-08-31: what runs next and why\secstatus{todo}}
\label{sec:launch}

The near-term strategy is one round on IGB: settle the objective question of
\cref{sec:campaign} Phase A with replicated arms at the incumbent's full budget, while the
incumbent's own checkpoint resumes past 9,900 epochs. Ten GPUs are available, six of
cabbi's eight RTX 6000 Ada (two left free for others) and all four mmli A100s, at two runs
per GPU, under the administration's 5-day-per-job limit, with no job chaining. Every design
choice below is stated with the measurement that forces it.

\subsection{Evidence 1: the replicate spread at 9,900 epochs is $0.0222$\secstatus{todo}}

Every expression run past 9,000 epochs on the leaderboard, eight of them, carries one
identical configuration: \file{quantile} head, lr $3 \times 10^{-4}$, dropout $0.1$,
$L = 6$, hidden 90, mask prior, \file{s1_pool}, seed 0. The script asserts this before
computing anything, so the spread below is run-to-run nondeterminism and nothing else.

\begin{table}[htbp]
\centering\small
%% SOURCE: results/launch_plan_evidence.json, replicate_spread_9900_epochs
\begin{tabular}{lrr}
\toprule
statistic & value & \\
\midrule
scores, sorted & \multicolumn{2}{l}{$0.1663,\ 0.1804,\ 0.1824,\ 0.1887,\ 0.2010,\ 0.2057,\ 0.2091,\ 0.2382$} \\
mean $\pm$ sd & $0.1965 \pm 0.0222$ & \\
expected maximum of 8 draws (Blom) & $0.2282$ & \\
observed maximum & $0.2382$ & at $z = 1.89$ \\
peak epochs, sorted & \multicolumn{2}{l}{$2199,\ 3398,\ 3503,\ 3715,\ 4109,\ 8859,\ 9691,\ 9697$} \\
\bottomrule
\end{tabular}
\caption[]{The eight identical-config expression runs at 9,900 epochs. The headline
$0.2382$ is the maximum of these draws; the config's central estimate is
$0.196 \pm 0.022$. Five of eight peaked by epoch 4,109 and never exceeded that value in the
following six thousand epochs, so budget raises the maximum draw rather than every draw.}
\label{tab:launch-replicates}
\end{table}

\subsection{Evidence 2: what that spread costs in replicates\secstatus{todo}}

\begin{table}[htbp]
\centering\small
%% SOURCE: results/launch_plan_evidence.json, power_80pct_alpha05 and
%% detectable_vs_n8_baseline
\begin{tabular}{rr@{\hspace{14mm}}rr}
\toprule
\multicolumn{2}{l}{two fresh arms} & \multicolumn{2}{l}{one fresh arm vs the $n = 8$ baseline} \\
gap to detect & runs per arm & runs in the arm & smallest detectable gap \\
\midrule
$0.02$ & $19.3$ & 2 & $0.049$ \\
$0.03$ & $8.6$  & 3 & $0.042$ \\
$0.04$ & $4.8$  & 6 & $0.034$ \\
$0.05$ & $3.1$  & 7 & $0.032$ \\
\bottomrule
\end{tabular}
\caption[]{Replicate arithmetic at $\sigma = 0.0222$, $80\,\%$ power, $\alpha = 0.05$. The
left block is the symmetric two-arm design; the right block is the design used here, where
each challenger head is compared against the eight quantile replicates that already exist.
The matched-budget head gaps on record ($+0.0085$ on the median,
\cref{tab:dist-confound}) sit below every row of this table, which is the arithmetic
demonstration that the previous round's one-run-per-arm scoring could not have resolved the
question it was asking.}
\label{tab:launch-power}
\end{table}

The consequence that shapes everything: \textbf{re-running \file{quantile} is wasted
budget}, because its $n = 8$ baseline already exists at this exact budget and config, and
every fresh slot spent on a challenger tightens the challenger side of an
already-anchored comparison.

\subsection{Evidence 3: wall clock, and the packing cost measured\secstatus{todo}}

From the same eight runs, the two IGB card classes separate cleanly on compute time and
not at all on wall clock: RTX 6000 Ada 17.2 s/epoch compute against the A100's 27.0, but
31.0 against 30.5 s/epoch wall, because validation and checkpointing overhead dominates.
So cabbi and mmli deliver the same epochs per day, and the partition split can be chosen on
availability alone. At those rates 9,900 epochs takes 3.5 days solo on either card, inside
the 5-day cap.

Packing was then measured directly rather than assumed, on two idle GilaHyper 46 GB cards:
the incumbent config resumed from the 1445 checkpoint with \file{train_eval_every=0}, once
solo and then two concurrent runs on one card, ten epochs each at steady state.
\src{experiments/019-simb-multimodal/scripts/gh\_packing\_benchmark.sh}

\begin{table}[htbp]
\centering\small
%% SOURCE: results/launch_plan_evidence.json, packing_benchmark
\begin{tabular}{lrr}
\toprule
 & solo & two packed on one GPU \\
\midrule
median s/epoch, per run & $17.0$ & $26.0$ \\
per-run slowdown & \multicolumn{2}{c}{$1.53\times$} \\
aggregate throughput & \multicolumn{2}{c}{$1.31\times$} \\
VRAM, both runs together & \multicolumn{2}{c}{$19{,}379$ MiB $= 9.5$ GiB per run} \\
\bottomrule
\end{tabular}
\caption[]{The packing benchmark. Memory is nowhere near binding, at $9.5$ GiB per run on
cards of 46 to 80 GiB; the whole cost of packing is the $1.53\times$ per-run slowdown.
Extrapolated to the IGB wall-clock rates, which is an extrapolation from a different
machine and an eval-free configuration and is labeled as such, a packed run covers
$\approx 9{,}100$ to $9{,}250$ epochs in 5 days, $7\,\%$ short of the 9,900 target.}
\label{tab:launch-packing}
\end{table}

Two runs per GPU is the chosen operating point: it doubles the replicate count, which
\cref{tab:launch-power} shows is the binding resource, at the price of arms reaching
$\approx 9{,}100$ rather than 9,900 epochs before the wall. The scoring rule below absorbs
that shortfall.

\subsection{The allocation: 20 slots\secstatus{todo}}

\begin{table}[htbp]
\centering\small
\begin{tabular}{llrlr}
\toprule
partition & GPUs & slots & arm & \file{max_epochs} \\
\midrule
mmli (A100) & 4 & 1 & resume \texttt{hx8pxdic} from its 9,900-epoch checkpoint & 19{,}000 \\
            &   & 7 & \file{dist=crps}, seeds 0--6 & 9{,}900 \\
cabbi (Ada) & 6 & 6 & \file{dist=laplace_crps}, seeds 0--5 & 9{,}900 \\
            &   & 6 & \file{dist=point}, seeds 0--5 & 9{,}900 \\
\bottomrule
\end{tabular}
\caption[]{The launch: 20 runs on 10 GPUs, two per GPU, everything else pinned to the
incumbent configuration. Nineteen fresh runs across three challenger heads, each scored
against the existing $n = 8$ \file{quantile} baseline, plus the resumed incumbent.}
\label{tab:launch-allocation}
\end{table}

\notesec{Justification, arm by arm}
\begin{description}
\item[Why these three heads and no others] \file{nll} is the designed $\sigma$-collapse
  negative control and is already last at every budget; \file{energy}'s covariance
  $\Sigma = D + VV^{\top}$ is global and does not enter $\mu$, so it cannot move a point
  metric (\cref{sec:campaign}); \file{quantile} needs no fresh runs because its baseline
  is the eight replicates of \cref{tab:launch-replicates}. That leaves exactly the three
  heads whose point estimates could differ: \file{crps} (Gaussian $\mu$), \file{laplace_crps}
  (Laplace median, structurally the closest analog to the quantile head's median knot), and
  \file{point} (all parameters on the location the metric reads, which is the corrected
  \texttt{O\_zmse} arm of \cref{sec:campaign}, since per-feature target z-scoring is
  already on).
\item[Why \file{crps} gets seven and the others six] It is the strongest challenger on the
  matched-budget record ($0.1631$ best, $0.1267$ median against \file{quantile}'s $0.1751$
  and $0.1352$), so its comparison deserves the tightest resolution: seven runs against
  the $n = 8$ baseline detects $0.032$; six detects $0.034$.
\item[Why full budget rather than 4,000 epochs] The comparison target is the incumbent at
  its own budget. Truncated arms would re-introduce the exact confound
  \cref{tab:dist-confound} documents, and the measured wall clock makes full budget
  affordable, so nothing is bought by truncating.
\item[Why the continuation, and why only one] The two best replicates peaked in their last
  $12\,\%$, so whether the top draw keeps climbing past 9,900 is open, and resuming the
  existing checkpoint is the only way to see past it without paying for the first 9,900
  again. It is one draw with $\sigma = 0.0222$, so it cannot establish saturation by
  itself; it can only show whether \emph{this} draw's curve turns. Budget: packed on an
  A100 at the extrapolated $46.7$ s/epoch, five days adds $\approx 9{,}200$ epochs, so
  \file{max_epochs} $= 19{,}000$. It runs on mmli, and it is a resume of the checkpoint
  rather than a chained job.
\item[Why seeds vary when the baseline's did not] The eight baseline runs spread $0.0222$
  at a \emph{fixed} seed label, so run-to-run nondeterminism alone produces the measured
  floor. Varying the seed across replicates adds initialization variance on top, making
  each arm's spread an honest upper-bound match to what any future user of the winning
  head would see.
\item[The scoring rule, declared before launch] Each arm and the baseline are scored as
  \file{roll_max} over epochs $\le 9{,}000$, the largest budget every packed run reaches
  before the 5-day wall, with the full-history \file{roll_max} reported alongside.
  \file{nmse}, the SD ratio $s$, and the coverage diagnostics are reported per arm as
  \cref{sec:campaign} requires, so a head that moves Pearson while leaving calibration
  broken is visible as such. No arm is compared on a raw per-epoch maximum.
\item[Operational constraints honored] 5 days per job per the administrator email, with a
  continuation counted as a fresh 5-day job rather than a chain; six of eight cabbi GPUs,
  leaving two free; all compute on partition nodes with the login node used only for
  \file{sbatch}, \file{rsync}, and \file{wandb}~\file{sync}.
\end{description}

\subsection{The model being launched: 1,194,509 parameters\secstatus{todo}}

The incumbent checkpoint's state dict reconciles exactly to the logged
\file{total_param_count} once two artifacts are removed: 98,370 parameters that appear
twice because \file{EquivariantPerturbationTransform} aliases layer 0 of its module lists
(\file{self.cross_attn = self.cross_attn_layers[0]} and siblings), and 12,273 non-trainable
buffers (the two per-gene normalization vectors and the 19 quantile knot positions).

\begin{table}[htbp]
\centering\small
%% SOURCE: results/launch_plan_evidence.json, param_count (1445-last.ckpt state dict,
%% deduplicated by tensor storage)
\begin{tabular}{lrl}
\toprule
module & parameters & role \\
\midrule
\file{transformer_layers} & 590{,}220 & 6 graph-masked encoder layers \\
\file{embedding_preprocessor} & 369{,}639 & input projection to $d = 90$ \\
\file{post_perturbation_mixing} & 101{,}430 & Perceiver mixing after the perturbation \\
\file{perturbation_transform} & 98{,}370 & the cross-attention perturbation operator \\
\file{perturbation_head} & 16{,}381 & graph-level scalar head \\
\file{per_gene_head} & 9{,}919 & per-gene readout, $d \to 19$ quantile knots \\
\file{observed_label_encoder} & 8{,}460 & masked teacher-forcing encoder \\
\file{cls_token} & 90 & \\
\midrule
total trainable & \textbf{1{,}194{,}509} & matches the logged \file{total_param_count} \\
\bottomrule
\end{tabular}
\caption[]{Parameter budget of the incumbent model. The striking row is the readout:
6,127 reporter genes are predicted by 9,919 parameters, because \file{PerGeneHead} is one
MLP shared across every gene token, emitting that gene's 19 quantile knots. The
equivariant weight sharing is what makes the output dimension nearly free, and it is also
why the perturbation operator, at $8\,\%$ of the parameters, carries the entire burden of
making predictions strain-specific: the encoder runs once on the wild-type graph, so
everything downstream of \file{perturbation_transform} is the only place strain identity
enters.}
\label{tab:launch-params}
\end{table}

\subsection{Graph attention, as currently built\secstatus{todo}}

For reading the arms above against the architecture, the current facts, all verified in
code or measured in \cref{sec:common}:

\begin{itemize}
\item The nine interaction graphs enter as a \textbf{hard additive attention mask} on the
  encoder's self-attention scores (\file{masked_fill} to $-\infty$ per head), not as edge
  features; composing $L = 6$ layers makes influence approximately 6-step random-walk
  reachability (\cref{sec:graph-probe}).
\item \file{_build_attention_mask} still \textbf{symmetrizes the two directed relations},
  discarding the measured TF $\to$ target signal ($0.5508$ against $0.5017$ rewired). The
  fix is queued in Phase B, not in this launch, which holds the architecture fixed by
  design.
\item The pair-form probe found the graph channel carries real signal on expression, up to
  AUC $0.6579$ on STRING database (\cref{tab:graph-pair-form}), so the masks stay.
\item The perturbation enters \emph{after} the encoder, through a cross-attention whose
  keys are the perturbed genes, and at a single deletion that softmax runs over one key
  and degenerates to $g(h_i + c_b)$: no term depends on the pair (deleted gene, reporter
  gene) (\cref{sec:expression}).
\end{itemize}

\subsection{Perturbation encodings that supply one-to-many, in this attention style
\secstatus{todo}}

The question the degeneracy raises: are there perturbation encodings, compatible with a
graph-masked transformer over gene tokens, in which \emph{one} perturbed gene produces
\emph{gene-specific} effects across all 6,127 reporters. Yes, on two tiers.

\notesec{Already implemented in \file{equivariant_cell_graph_transformer.py}, each gated to
be the exact identity at initialization} Ordered by the rank of the pair term they add:

\begin{table}[htbp]
\centering\small
%% SOURCE: torchcell/models/equivariant_cell_graph_transformer.py (mechanisms and ranks
%% from the module docstrings and init code; the basis64 score from round_leaderboards.csv)
\begin{tabular}{llp{62mm}}
\toprule
mechanism & pair rank & how one deletion reaches many genes \\
\midrule
additive reference & 0 & it does not: one $c_b$ added to every gene \\
null sink & 9 & sigmoid gate per head, $\alpha_i$ depends on the querying gene \\
graph propagation & data & random-walk mass from the deleted gene over the nine graphs, plus a hop-0 self indicator \\
low-rank bilinear & $r$ & $\sum_j b_{ij}(h_i)\, a_j(c_b)$ \\
response basis (\texttt{P\_basis64}) & 64 & measured at $0.2274$ in 2.0 days \\
Hadamard / FiLM & 90 & $h_i \odot (1 + \gamma(c_b))$, gene and strain multiply \\
\bottomrule
\end{tabular}
\caption[]{In-repo one-to-many perturbation mechanisms, all resident in the model file and
selectable by config. These are the Phase B arms of \cref{tab:phase-b}; none is in this
launch, which settles the objective first so mechanism deltas are measured under a settled
loss.}
\label{tab:launch-pair-mechanisms}
\end{table}

\notesec{The published designs, mapped onto this architecture} These are literature
descriptions, not measurements made here.

\begin{itemize}
\item \textbf{GEARS} (Roohani et al., 2024) propagates a learned perturbation embedding
  over a relation graph to produce per-gene deltas. \file{PerturbationGraphPropagation} is
  this mechanism with fixed random-walk features in place of learned message passing, so
  the cheap version is already an arm; the learned version is the upgrade if the cheap one
  moves anything.
\item \textbf{Token-edit conditioning} (Geneformer's in-silico deletion, Theodoris et al.
  2023; scGPT's perturbation flag, Cui et al. 2024; STATE-style perturbation tokens) puts
  the perturbation \emph{into the encoder's input}, as removal of the gene's token or a
  condition embedding added to it, so the encoder's own self-attention, here graph-masked,
  carries the deletion outward and the one-key degeneracy never arises: every gene attends
  to the perturbed gene among 6,607 tokens, not to a key set of size one. This is the
  design that most directly ``fits into our attention style''. Its cost is structural: the
  current model runs the encoder \textbf{once} on the wild-type graph and reuses it for
  every strain, which is precisely why the perturbation had to enter post hoc; a token
  edit forces one encoder pass per strain, multiplying epoch cost by roughly the batch
  size. \textbf{Hypothesis, untested:} at 39 batches per epoch and 17 to 31 s/epoch, that
  overhead is a wall-clock multiplier this campaign cannot absorb, but a 2-layer
  strain-specific encoder atop frozen wild-type layer outputs could cap it.
\item \textbf{CPA-style additive latent composition} (Lotfollahi et al., 2023) is the same
  rank-0 additive structure as the reference operator, with gene specificity living
  entirely in the decoder. It is what the model already is, and its limitation on the pair
  axis is the one \cref{sec:expression} proves.
\end{itemize}

The ordering stands as in \cref{sec:campaign}: objective first, with this launch; the
mechanism tier second, under the winning objective, with the in-repo arms before any
encoder restructuring.
